\rtmtight
\begin{tblr}{width=\textwidth,colspec={|Q[c,m,2.1cm]|X[l,m]|},hlines,vlines,hspan=minimal,rowsep=1.5pt,colsep=3pt,row{1}={bg=headgray,font=\bfseries}}
\SetCell{c,m}\hfil\logo\hfil & \textbf{POLITEKNIK NEGERI MALANG}\par\textbf{JURUSAN TEKNOLOGI INFORMASI}\par\textbf{PROGRAM STUDI: D4 TEKNIK INFORMATIKA} \\
\SetCell[c=2]{l} \textbf{RENCANA TUGAS MAHASISWA (RTM) DAN RUBRIK PENILAIAN} & \\
Nama Mata Kuliah & Kewirausahaan Berbasis Teknologi \\
Kode Mata Kuliah & RTI255001 \quad \textbf{sks / Jam perminggu:} 2 SKS/4 jam \quad \textbf{Semester:} 5 \\
Dosen Pengampu & Tim Dosen Mata Kuliah \\
\end{tblr}
\assessmentblock{Business Game dan Aktivitas Kewirausahaan}{Case Method dan simulasi bisnis}{SCPMK0101-03701, SCPMK0301-03703}{Mahasiswa berpartisipasi aktif dalam simulasi bisnis (Business Game), diskusi kasus kewirausahaan, dan workshop Business Model Canvas untuk mengembangkan jiwa dan kompetensi wirausaha.}{Mengikuti simulasi bisnis, berdiskusi kelompok, menganalisis kasus, dan mengerjakan modul Business Game secara aktif.}{Modul Business Game yang terisi, catatan refleksi, dan hasil diskusi kelompok.}{Partisipasi aktif dan kontribusi dalam simulasi bisnis & 10 & Mahasiswa berpartisipasi penuh di setiap sesi, memberikan kontribusi bermakna dalam diskusi kelompok, dan menyelesaikan semua modul Business Game. & Mahasiswa hadir dan berpartisipasi di sebagian besar sesi, tetapi kontribusi dalam diskusi kelompok belum konsisten. & Partisipasi sangat terbatas, modul Business Game tidak diselesaikan, atau tidak ada kontribusi bermakna. \\ \\
Kualitas pengambilan keputusan dalam Business Game & 8 & Keputusan bisnis dalam simulasi diambil berdasarkan analisis data yang tersedia, alasan keputusan logis, dan hasilnya direfleksikan secara kritis. & Keputusan diambil dengan pertimbangan cukup, tetapi analisis kurang mendalam atau refleksi terbatas. & Keputusan diambil tanpa dasar analisis atau tidak ada refleksi terhadap hasil simulasi. \\ \\
Keterlibatan dalam workshop Business Model Canvas & 7 & BMC dikerjakan lengkap 9 komponen dengan analisis yang kontekstual dan diskusi aktif. & BMC dikerjakan sebagian besar komponen, tetapi analisis atau diskusi kurang mendalam. & BMC tidak diselesaikan atau diisi tanpa pemahaman yang memadai. \\}

\assessmentblock{Tugas Rencana Bisnis Teknologi}{Penugasan terstruktur}{SCPMK0101-03701, SCPMK0101-03702, SCPMK0301-03703}{Mahasiswa mengerjakan serangkaian tugas terstruktur yang membangun rencana bisnis teknologi secara bertahap, mulai dari ide produk, analisis pasar, aspek keuangan, hingga BMC lengkap.}{Mengerjakan tugas individual atau kelompok yang diberikan setiap pertemuan, mengumpulkan pada tenggat yang ditentukan.}{Dokumen tugas terstruktur, laporan analisis pasar, lembar kerja keuangan, dan BMC.}{Kedalaman analisis pasar dan kelayakan produk & 12 & Analisis pasar mencakup segmen target, ukuran pasar, kompetitor, dan permintaan dengan data pendukung yang relevan. & Analisis pasar mencakup elemen utama, tetapi data pendukung terbatas atau analisis kompetitor kurang mendalam. & Analisis pasar sangat dangkal, tidak berbasis data, atau tidak menunjukkan pemahaman target pasar. \\ \\
Kelengkapan aspek teknis, keuangan, dan BMC & 13 & Aspek teknis produk, proyeksi keuangan (laba rugi, BEP, ROI), dan 9 komponen BMC diselesaikan dengan lengkap dan konsisten satu sama lain. & Sebagian besar aspek dikerjakan, tetapi beberapa komponen keuangan atau BMC masih kurang lengkap atau belum konsisten. & Banyak komponen tidak dikerjakan atau terdapat inkonsistensi yang signifikan antar bagian rencana bisnis. \\ \\
Kualitas dokumentasi dan laporan bisnis & 10 & Laporan terstruktur rapi, menggunakan bahasa yang tepat, data disajikan dalam tabel/grafik yang informatif, dan kesimpulan logis. & Laporan cukup terstruktur, tetapi beberapa bagian kurang rapi atau penyajian data kurang optimal. & Laporan tidak terstruktur, sulit dibaca, atau tidak mencerminkan pemahaman terhadap materi. \\}

\assessmentblock{Kuis Kewirausahaan}{Kuis tertulis}{SCPMK0101-03701, SCPMK0101-03702}{Mahasiswa menjawab soal pilihan ganda dan uraian singkat terkait konsep kewirausahaan, karakteristik wirausaha, analisis pasar, dan aspek teknis bisnis.}{Mengerjakan soal secara mandiri dalam waktu yang ditentukan.}{Lembar jawaban kuis.}{Pemahaman konsep kewirausahaan dan karakteristik wirausaha & 5 & Konsep kewirausahaan, tipe wirausaha, intensi kewirausahaan, dan karakteristik PEC dijelaskan dengan tepat dan lengkap. & Sebagian besar konsep tepat, tetapi beberapa aspek masih kurang akurat atau tidak lengkap. & Jawaban salah atau hanya memberikan definisi tanpa pemahaman. \\ \\
Ketepatan analisis pasar dan aspek bisnis & 5 & Pertanyaan tentang analisis pasar, permintaan, dan aspek teknis bisnis dijawab dengan tepat berdasarkan pemahaman materi. & Sebagian besar pertanyaan dijawab dengan benar, tetapi beberapa kurang akurat. & Banyak pertanyaan dijawab salah atau jawaban tidak relevan dengan materi. \\}

\assessmentblock{UTS Aspek Teknis dan Keuangan Bisnis}{Ujian Tengah Semester}{SCPMK0101-03701, SCPMK0101-03702}{Mahasiswa mengerjakan soal UTS yang mencakup konsep kewirausahaan, analisis pasar, strategi pemasaran, dan perhitungan keuangan dasar bisnis.}{Mengerjakan soal secara mandiri.}{Lembar jawaban UTS.}{Pemahaman analisis pasar dan strategi pemasaran & 4 & Mahasiswa dapat menganalisis segmen pasar, menghitung permintaan, dan menyusun strategi pemasaran 4P dengan tepat untuk kasus yang diberikan. & Analisis pasar dan strategi pemasaran cukup tepat, tetapi beberapa aspek kurang mendalam. & Tidak dapat menganalisis pasar atau strategi pemasaran tidak relevan dengan kasus. \\ \\
Kemampuan menyusun laporan keuangan dasar & 3 & Laporan laba rugi sederhana, buku kas, atau perhitungan BEP disusun dengan benar dan logis. & Struktur laporan keuangan sebagian besar benar, tetapi terdapat kesalahan pada satu atau dua perhitungan. & Laporan keuangan tidak dapat disusun atau terdapat kesalahan fundamental. \\ \\
Ketepatan analisis aspek teknis bisnis & 3 & Aspek teknis produk/layanan teknologi diidentifikasi dengan tepat, kapasitas produksi atau kebutuhan infrastruktur diestimasi logis. & Aspek teknis diidentifikasi sebagian besar benar, tetapi estimasi kurang akurat. & Aspek teknis tidak teridentifikasi atau tidak relevan dengan jenis bisnis yang dibahas. \\}

\assessmentblock{UAS Presentasi Rencana Bisnis Teknologi}{UAS presentasi dan pitch}{SCPMK0301-03704}{Mahasiswa mempresentasikan rencana bisnis teknologi lengkap yang telah disusun selama satu semester di hadapan penguji, mempertahankan ide bisnis, dan menjawab pertanyaan.}{Mempresentasikan rencana bisnis, menjawab pertanyaan penguji, dan menunjukkan kesiapan bisnis.}{Dokumen rencana bisnis final, slide presentasi, dan/atau poster bisnis.}{Kelengkapan dan kohesi rencana bisnis teknologi & 8 & Rencana bisnis mencakup semua komponen (latar belakang, analisis pasar, BMC, aspek teknis, keuangan, hukum, SDM) yang saling terhubung dan konsisten. & Sebagian besar komponen tersedia, tetapi beberapa bagian belum lengkap atau koneksi antar komponen kurang jelas. & Rencana bisnis tidak lengkap, banyak komponen hilang, atau komponen tidak konsisten satu sama lain. \\ \\
Kualitas presentasi dan kemampuan pitching & 7 & Presentasi terstruktur dengan baik, pesan bisnis disampaikan dengan jelas dan meyakinkan, slide informatif, dan penampilan profesional. & Presentasi cukup baik, pesan bisnis tersampaikan, tetapi ada beberapa bagian yang kurang jelas atau slide kurang informatif. & Presentasi tidak terstruktur, pesan bisnis tidak tersampaikan dengan jelas, atau penampilan tidak profesional. \\ \\
Kemampuan menjawab pertanyaan dan argumentasi & 5 & Semua pertanyaan dijawab dengan tepat, argumen didukung data dan analisis yang relevan, serta menunjukkan penguasaan materi yang mendalam. & Sebagian besar pertanyaan dijawab dengan benar, tetapi beberapa argumen kurang kuat atau tidak didukung data. & Banyak pertanyaan tidak dapat dijawab atau argumen tidak relevan dengan pertanyaan yang diajukan. \\}

\begin{tblr}{width=\textwidth,colspec={|Q[l,m,3.2cm]|X[l,m]|},hlines,vlines,hspan=minimal,row{1-3}={bg=headgray},rowsep=1.5pt,colsep=3pt}
Jadwal Pelaksanaan: & Minggu 1--17 sesuai RPS. Setiap evaluasi memiliki RTM dan rubrik multi-kriteria. \\
Lain-Lain Yang Diperlukan: & LMS, modul business game, perangkat lunak Microsoft Office, Canva/Figma (opsional), dan perangkat presentasi. \\
Pustaka: & Mengacu pustaka utama dan pendukung pada bagian RPS. \\
\end{tblr}
